\chapter{Rehabilitation Medicine}

\medterm{Adaptive Equipment}-- Devices that help individuals with disabilities perform daily activities independently. Examples: reachers, dressing sticks, button hooks, ergonomic utensils, grab bars, shower chairs, wheelchairs, walkers, canes. Selection based on individual needs and goals. Occupational therapy assessment recommended.

\medterm{Adaptive Sports}-- Sports modified for participants with disabilities. Examples: wheelchair basketball, para swimming, adaptive skiing, sitting volleyball. Promotes physical fitness, psychosocial well-being, inclusion. Paralympic Games showcase elite adaptive athletes.

\medterm{Amputation}-- Surgical removal of body part. Types: partial/complete, above/below knee, transhumeral/transradial. Indications: trauma, vascular disease, malignancy, infection. Rehabilitation: prosthetic fitting, gait training, strength, ADL independence, psychological support. Phantom limb pain common.

\medterm{Assistive Devices}-- Tools enhancing mobility and independence. Categories: mobility (walkers, crutches, canes, wheelchairs), ADL (reachers, dressing aids, adaptive utensils), communication (sign language, communication boards), environmental (retractable showerheads, raised toilet seats). Prescription by PT/OT.

\medterm{Balance Training}-- Exercises improving postural stability and preventing falls. Components: static balance, dynamic balance, reactive balance. Techniques: weight shifting, tandem stance, single-leg stance, perturbations, balance boards. Essential for elderly, neurological conditions.

\medterm{Cardiac Rehabilitation}-- Supervised program for cardiac patients post-MI, PCI, CABG, heart failure. Phases: inpatient (early mobilization), outpatient (supervised exercise, education), maintenance. Components: exercise training, risk factor modification, psychological support, return to work counseling. Reduces mortality, hospitalizations.

\medterm{Cerebral Palsy}-- Non-progressive motor disorder from brain injury before, during, or after birth. Types: spastic (most common), dyskinetic, ataxic, mixed. Symptoms: impaired movement, posture, muscle tone. Associated: intellectual disability, epilepsy, speech/language disorders. Management: multidisciplinary—PT, OT, speech therapy, medications, orthotics, surgery.

\medterm{Chronic Fatigue Syndrome}-- See myalgic encephalomyelitis.

\medterm{Concussion}-- Mild traumatic brain injury causing transient neurological dysfunction. Symptoms: headache, confusion, memory problems, dizziness, nausea, sensitivity to light/noise. Management: physical and cognitive rest, gradual return to activity, symptom-limited exercise. Most recover within 7-10 days.

\medterm{COPD Rehabilitation}-- Comprehensive intervention for COPD patients including exercise training, education, behavior change. Components: aerobic exercise, strength training, breathing techniques, nutrition, psychosocial support. Improves exercise capacity, quality of life, reduces hospitalizations.

\medterm{Dysphagia Rehabilitation}-- See rehabilitation.

\medterm{End-Stage Renal Disease}-- See nephrology.

\medterm{Fall Prevention}-- See geriatrics.

\medterm{Gait Training}-- Teaching or retraining walking pattern. Techniques: parallel bars, walker, cane, stair training. Addresses abnormalities: limping, antalgic gait, hemiplegic gait, parkinsonian gait. Goals: safety, efficiency, independence.

\medterm{Hydrotherapy}-- Exercise in water using buoyancy, resistance, warmth. Benefits: reduced weight-bearing, improved circulation, pain relief, increased range of motion. Used for arthritis, spinal cord injury, amputation, neurological conditions.

\medterm{Isokinetic Exercise}-- Muscle contraction at constant speed throughout range of motion. Uses specialized equipment. Assesses and improves strength.

\medterm{Joint Replacement Rehabilitation}-- See orthopedics.

\medterm{Kinesio Taping}-- See physical therapy.

\medterm{Lymphedema Management}-- See geriatrics.

\medterm{Myalgic Encephalomyelitis/Chronic Fatigue Syndrome}-- ME/CFS: debilitating condition characterized by profound fatigue, post-exertional malaise, unrefreshing sleep, cognitive impairment. Diagnosis: exclusion. Treatment: pacing, symptom management, graded exercise therapy controversial.

\medterm{Neurological Rehabilitation}-- Rehabilitation for stroke, spinal cord injury, traumatic brain injury, Parkinson disease, multiple sclerosis. Goals: maximize function, independence, quality of life. Multidisciplinary team: physiatrist, PT, OT, speech therapy, psychology.

\medterm{Occupational Therapy}-- OT helps patients perform daily activities (ADLs, IADLs) through therapeutic activities, adaptive equipment, environmental modification. Focus on upper extremity function, cognition, sensory integration.

\medterm{Pain Management}-- See emergency medicine.

\medterm{Post-Surgical Rehabilitation}-- Recovery program after surgery. Phases: acute (pain control, wound care), subacute (range of motion, strengthening), functional (return to activities). Protocols vary by procedure.

\medterm{Prosthesis}-- Artificial replacement for missing body part. Types: upper extremity, lower extremity. Components: socket, suspension, joint, pylon, terminal device. Fitting: certified prosthetist. Training: physical therapy.

\medterm{Psychological Support}-- Counseling, therapy for patients with chronic illness, disability, pain. Address depression, anxiety, adjustment difficulties, coping strategies.

\medterm{Quality of Life}-- See geriatrics.

\medterm{Respiratory Rehabilitation}-- See pulmonology.

\medterm{Return to Work}-- See occupational medicine.

\medterm{Spinal Cord Injury}-- See neurology.

\medterm{Therapeutic Exercise}-- Exercise prescribed to improve function, reduce pain, prevent disability. Types: stretching, strengthening, aerobic, balance, functional. Progression: low to high intensity.

\medterm{Transitional Care}-- See geriatrics.

\medterm{Urinary Retraining}-- See urology.

\medterm{Vestibular Rehabilitation}-- See neurology.

\medterm{Wheelchair}-- Mobility device for patients with walking difficulties. Types: manual, power. Components: seat, backrest, wheels, footrests, armrests. Assessment: seating evaluation, pressure relief, posture. Training: mobility, transfers, safety.

\medterm{Acute Rehabilitation}-- Intensive inpatient rehabilitation (3+ hours therapy/day) for acute conditions: stroke, spinal cord injury, traumatic brain injury, major trauma. Length of stay: 1-4 weeks.

\medterm{Inpatient Rehabilitation}-- See acute rehabilitation.

\medterm{Outpatient Rehabilitation}-- Therapy sessions 2-3 times per week. For conditions requiring ongoing therapy: post-surgical, chronic pain, neurological conditions.

\medterm{Home Health Rehabilitation}-- Therapy provided at home. For patients unable to travel. Home exercise program emphasized.

\medterm{Subacute Rehabilitation}-- Less intensive than acute rehab. For patients needing continued therapy but not as intensive. Skilled nursing facility setting.

\medterm{Long-Term Rehabilitation}-- Chronic condition management. Focus on maintaining function, preventing complications.

\medterm{Intensive Care Rehabilitation}-- Early rehabilitation in ICU setting. Early mobilization, positioning, passive range of motion. Improves outcomes, reduces complications.

\medterm{Multidisciplinary Team}-- Rehabilitation team includes: physiatrist, physical therapist, occupational therapist, speech-language pathologist, psychologist, social worker, case manager, nursing, recreation therapist.

\medterm{ICF Framework}-- International Classification of Functioning, Disability and Health. Biopsychosocial model: body functions/structures, activities, participation, environmental factors, personal factors.

\medterm{Functional Independence Measure}-- FIM: standardized assessment of disability. Scores 18 (total assistance) to 7 (independent). 18 items: self-care, sphincter control, transfers, locomotion, communication, social cognition.

\medterm{Barthel Index}-- See geriatrics.

\medterm{Katz Index}-- Assessment of ADLs: bathing, dressing, toileting, transferring, continence, feeding.

\medterm{Lawton IADL Scale}-- Instrumental ADLs: phone use, shopping, food preparation, housekeeping, laundry, transportation, medication management, finance management.

\medterm{Motivation}-- Key factor in rehabilitation success. Intrinsic vs extrinsic motivation. Motivational interviewing technique.

\medterm{Goal Setting}-- SMART goals: Specific, Measurable, Achievable, Relevant, Time-bound. Patient-centered goals improve adherence.

\medterm{Outcome Measurement}-- Standardized tools: FIM, Barthel, MRC, ASIA (spinal cord), NIHSS (stroke), etc. Track progress, adjust treatment.

\medterm{Evidence-Based Practice}-- Integration of best research evidence, clinical expertise, patient values. Rehabilitation guidelines based on systematic reviews.

\medterm{Neuroplasticity}-- Brain's ability to reorganize by forming new neural connections. Basis for rehabilitation after neurological injury. Use-dependent plasticity: repetitive practice drives recovery.

\medterm{Constraint-Induced Movement Therapy}-- CIMT: restrain unaffected limb, force use of affected limb. Based on neuroplasticity. Effective for stroke hemiparesis.

\medterm{Mirror Therapy}-- Mirror reflection of unaffected limb movement creates visual illusion of affected limb moving. Reduces phantom limb pain, improves stroke recovery.

\medterm{Robot-Assisted Rehabilitation}-- Robotic devices for repetitive task training. Upper extremity, lower extremity, gait training. Objective measurement, high repetition.

\medterm{Virtual Reality}-- Immersive technology for rehabilitation. Engaging, motivates exercise. Used for balance, coordination, cognitive training.

\medterm{Tele-Rehabilitation}-- Remote rehabilitation via technology. Video consultations, remote monitoring, home exercise programs. Increases access, especially rural.

\medterm{Patient Education}-- Teach patient and family about condition, exercises, equipment, prevention. Improves adherence, outcomes.

\medterm{Family Involvement}-- Family support crucial for rehabilitation success. Training family as therapists, addressing caregiver burden.

\medterm{Community Reintegration}-- Return to community roles: work, leisure, social. Vocational rehabilitation, accessibility modifications.

\medterm{Vocational Rehabilitation}-- Help patients return to work. Job analysis, accommodations, skills training, job placement.

\medterm{Disability Evaluation}-- Assessment of functional limitations for legal, insurance, or benefits purposes. Objective measures, standardized tests.

\medterm{Palliative Rehabilitation}-- Rehabilitation for patients with terminal illness. Focus on comfort, function, quality of life.

\medterm{Hospice Rehabilitation}-- Rehabilitation in hospice setting. Symptom management, positioning, caregiver education.

\medterm{Pediatric Rehabilitation}-- Rehabilitation for children with congenital or acquired disabilities. Developmental considerations, family-centered care, school collaboration.

\medterm{Geriatric Rehabilitation}-- See geriatrics.

\medterm{Sports Rehabilitation}-- See sports medicine.

\medterm{Occupational Rehabilitation}-- See occupational medicine.

\medterm{Cardiac Rehabilitation}-- See cardiac rehabilitation above.

\medterm{Pulmonary Rehabilitation}-- See pulmonology.

\medterm{Neurological Rehabilitation}-- See neurological rehabilitation above.

\medterm{Orthopedic Rehabilitation}-- See orthopedics.

\medterm{Oncology Rehabilitation}-- Cancer rehabilitation: fatigue, pain, neuropathy, lymphedema, functional impairment. Pre- and post-treatment support.

\medterm{Integrative Rehabilitation}-- Combination of conventional and complementary therapies: acupuncture, massage, yoga, meditation.

\medterm{Rehabilitation Engineering}-- Technology development for disability: prosthetics, orthotics, assistive devices, environmental control.

\medterm{Biomechanics}-- Study of movement mechanics. Gait analysis, joint forces, muscle activation. Applied to rehabilitation device design.

\medterm{Electromyography}-- EMG: measurement of muscle electrical activity. Used for diagnosis, biofeedback, device control.

\medterm{Motion Analysis}-- Assessment of movement patterns. 3D gait analysis, video analysis. Identifies abnormalities, guides treatment.

\medterm{Force Plates}-- Measure ground reaction forces during gait. Assess balance, symmetry, loading.

\medterm{Balance Assessment}-- Tests: Berg Balance Scale, Tinetti, timed up and go, Romberg. Identifies fall risk.

\medterm{Strength Testing}-- Manual muscle testing, dynamometry, isokinetic testing. Quantifies weakness, tracks progress.

\medterm{Range of Motion}-- ROM: joint movement capacity. Goniometer measurement. Restrictions from contracture, pain, swelling.

\medterm{Flexibility}-- Tissue extensibility. Stretching programs improve flexibility.

\medterm{Endurance}-- Aerobic capacity, muscular endurance. Cardiorespiratory fitness testing.

\medterm{Coordination}-- Smooth, efficient movement. Tests: finger-to-nose, heel-to-shin, rapid alternating movements.

\medterm{Agility}-- Ability to change direction quickly. Essential for sports, functional mobility.

\medterm{Power}-- Force x velocity. Important for functional tasks, fall prevention.

\medterm{Reaction Time}-- Time from stimulus to response. Declines with age, neurological disease.

\medterm{Walking Speed}-- Gait speed: important functional marker. <0.8 m/s indicates disability risk.

\medterm{Transfer Training}-- Teaching safe transfer techniques: bed to chair, chair to toilet, floor to standing. Use of transfer boards, lift equipment.

\medterm{Wheelchair Skills}-- Training for wheelchair users: propulsion, maneuvering, obstacles, slopes, stairs, transfers.

\medterm{Prosthetic Training}-- Training for amputees: stump care, donning/doffing, gait training, balance, stair negotiation, community ambulation.

\medterm{Orthotic Training}-- Training for brace users: donning/doffing, skin care, gait training with brace.

\medterm{Assistive Device Training}-- Training for walkers, canes, crutches: fitting, gait patterns, stair negotiation, safety.

\medterm{ADL Training}-- Training for activities of daily living: dressing, bathing, grooming, feeding, toileting. Adaptive techniques, equipment.

\medterm{IADL Training}-- Training for instrumental ADLs: cooking, cleaning, shopping, managing medications, finances.

\medterm{Pain Management}-- Multimodal pain management: medications, physical modalities, cognitive-behavioral therapy, intervention.

\medterm{Spasticity Management}-- See neurology.

\medterm{Contracture Prevention}-- Positioning, stretching, splinting, orthotics, serial casting, surgery.

\medterm{Pressure Ulcer Prevention}-- See geriatrics.

\medterm{Bladder Retraining}-- See urology.

\medterm{Bowel Program}-- See gastroenterology.

\medterm{Sexual Rehabilitation}-- Address sexual function after injury, illness. Counseling, devices, medications.

\medterm{Vocational Assessment}-- Evaluate work capacity, skills, interests. Match to job opportunities.

\medterm{Driving Evaluation}-- Assess driving fitness after neurological, musculoskeletal conditions. On-road testing, vehicle modifications.

\medterm{Driving Rehabilitation}-- Training for adaptive driving: hand controls, left-foot accelerator, other modifications.

\medterm{Swallowing Therapy}-- See speech-language pathology.

\medterm{Speech Therapy}-- See speech-language pathology.

\medterm{Language Therapy}-- See speech-language pathology.

\medterm{Cognitive Rehabilitation}-- See neurology.

\medterm{Behavioral Rehabilitation}-- Behavioral interventions for chronic conditions: pain management, addiction, adjustment.

\medterm{Psychosocial Rehabilitation}-- Support for mental illness recovery: skills training, supported employment, peer support.

\medterm{Recreational Therapy}-- Use of recreation for therapeutic purposes. Improves quality of life, socialization, coping.

\medterm{Music Therapy}-- Use of music for therapeutic purposes. Pain reduction, emotional expression, cognitive stimulation.

\medterm{Art Therapy}-- Use of art-making for therapeutic purposes. Emotional expression, coping, self-awareness.

\medterm{Animal-Assisted Therapy}-- Use of animals for therapeutic purposes. Reduces stress, improves mood, motivates exercise.

\medterm{Hippotherapy}-- Therapeutic horseback riding for neurological, musculoskeletal conditions. Improves balance, posture, coordination.

\medterm{Aquatic Therapy}-- See hydrotherapy.

\medterm{Thermal Therapy}-- Heat and cold therapy. Heat: relax muscles, increase circulation. Cold: reduce inflammation, pain.

\medterm{Electrical Stimulation}-- TENS, NMES for pain, muscle re-education, spasticity.

\medterm{Ultrasound Therapy}-- Therapeutic ultrasound for tissue healing, pain, scar tissue.

\medterm{Light Therapy}-- Phototherapy for skin conditions, seasonal affective disorder, jaundice.

\medterm{Magnetic Therapy}-- PEMF for bone healing, pain. Evidence mixed.

\medterm{Hydrotherapy}-- See hydrotherapy above.

\medterm{Kinesiology}-- Study of human movement. Applied to rehabilitation, sports, ergonomics.

\medterm{Biomechanics}-- See rehabilitation engineering.

\medterm{Ergonomics}-- Design for human use. Workstation setup, tool design, injury prevention.

\medterm{Workplace Assessment}-- Evaluate job demands, worker capabilities, modify job or environment.

\medterm{Job Modification}-- Adjust job duties, equipment, schedule for disability.

\medterm{Reasonable Accommodation}-- See occupational medicine.

\medterm{Disability Insurance}-- See occupational medicine.

\medterm{Workers Compensation}-- See occupational medicine.

\medterm{Return-to-Work}-- Gradual return to work after illness, injury. Phased return, modified duties.

\medterm{Stay-at-Work}-- Prevent work disability through early intervention, accommodation.

\medterm{Occupational Therapy Assessment}-- See occupational therapy.

\medterm{Physical Therapy Assessment}-- Comprehensive evaluation: history, observation, ROM, strength, sensation, coordination, balance, gait, functional tasks.

\medterm{Speech-Language Pathology Assessment}-- Evaluate communication, swallowing, cognition. Standardized tests, clinical observation.

\medterm{Rehabilitation Medicine}-- Physician specialty focusing on function, quality of life for patients with impairment, disability. Also called physiatry.

\medterm{Physiatrist}-- See rehabilitation medicine.

\medterm{Rehabilitation Hospital}-- Inpatient facility specializing in rehabilitation.

\medterm{Rehabilitation Unit}-- Hospital ward dedicated to rehabilitation.

\medterm{Rehabilitation Clinic}-- Outpatient rehabilitation facility.

\medterm{Home Rehabilitation}-- See home health rehabilitation.

\medterm{Community Rehabilitation}-- Rehabilitation services in community settings.

\medterm{School-Based Rehabilitation}-- Rehabilitation services in schools for children.

\medterm{Early Intervention}-- Rehabilitation for infants, toddlers with developmental delays or disabilities.

\medterm{Transition to Adulthood}-- Support for youth with disabilities transitioning to adult services.

\medterm{Care Coordination}-- Coordination of rehabilitation services across settings, providers.

\medterm{Case Management}-- See geriatrics.

\medterm{Discharge Planning}-- Planning for post-hospitalization care. Early identification of needs, arrangements.

\medterm{Follow-Up}-- Post-rehabilitation follow-up. Monitor progress, adjust treatment, prevent complications.

\medterm{Outcome Research}-- Study of rehabilitation outcomes. Effectiveness, cost-effectiveness, patient-reported outcomes.

\medterm{Quality Improvement}-- Continuous improvement of rehabilitation services.

\medterm{Patient-Centered Care}-- Care respecting patient preferences, needs, values. Shared decision-making.

\medterm{Cultural Competence}-- Care respectful of cultural differences. Language services, cultural adaptation.

\medterm{Health Literacy}-- Patient ability to understand health information. Teach-back method, plain language.

\medterm{Self-Management}-- Patient managing own condition. Education, skills, support.

\medterm{Self-Efficacy}-- Patient belief in ability to manage condition. Influences adherence, outcomes.

\medterm{Adherence}-- Patient following treatment plan. Barriers: cost, side effects, complexity, beliefs. Strategies: simplification, education, motivational interviewing.

\medterm{Therapeutic Alliance}-- Collaboration between therapist and patient. Foundation for successful rehabilitation.
